\subsection{Architectural Result Priority}
When one instruction could encounter more than one exceptional condition, the architecturally reported result is
the first applicable condition in the following order.

\begin{enumerate}
\item Instruction fetch and fetch translation.
\item Instruction framing and acquisition of the complete instruction record.
\item Opcode recognition.
\item Availability of every extension required by the decoded operation.
\item Fixed-field, reserved-field, and effective-address-form legality.
\item Privilege checks.
\item Selector, control-image, and other architectural-state validation.
\item Predicate evaluation for a conditional operation.
\item Every enabled operand access in architectural access order.
\item Execution-stage faults defined by the decoded operation.
\end{enumerate}

A false predicate suppresses the conditional operation\textquotesingle{}s target-memory and implicit stack accesses. It does not
suppress instruction framing, (:term:base.terms.opcode:) and extension checks, encoding legality, privilege checks, or control-state
validation. Explicit operands are accessed in instruction operand order unless an instruction states another order.
For a memory-to-memory operation, the complete source read precedes access to the destination, so a source-access
failure has priority over a destination-access failure.

Implicit accesses use the order stated by the instruction. In particular, PUSH captures its source before accessing
the destination stack slot; POP reads the stack slot before validating or writing its destination; PUSHP uses listed
register order and POPP uses reverse listed order; CALL validates its target before accessing the return stack; RET
reads the return stack before validating its target; and (:term:base.terms.long_call:), (:term:base.terms.long_return:), Event Return, and supervisor-entry
instructions use the step order in their instruction descriptions.
